\textbf{Phần 1. Trắc nghiệm lựa chọn}
\setcounterpageref{ex}{0}

\begin{ex} %[3H3N16]
Ba điểm $A, M, B$ thẳng hàng theo thứ tự đó. Điểm nằm ở giữa hai điểm $A$ và $B$ là
\choice
{Điểm $A$}
{\True Điểm $M$}
{Điểm $B$}
{Không có điểm nào}
\loigiai{
Điểm $M$ nằm giữa hai điểm $A$ và $B$.\\
\textbf{Đáp án của câu hỏi là B}
}
\end{ex}

\begin{ex} %[3H3N16]
Điểm $M$ được gọi là trung điểm của đoạn thẳng $AB$ khi
\choice
{\True $M$ nằm giữa $A, B$ và $AM = MB$}
{$M$ nằm ngoài đoạn thẳng $AB$}
{$M$ trùng với điểm $A$}
{$M$ trùng với điểm $B$}
\loigiai{
Trung điểm là điểm nằm giữa và chia đoạn thẳng thành hai phần bằng nhau: $AM = MB$.\\
\textbf{Đáp án của câu hỏi là A}
}
\end{ex}

\begin{ex} %[3H3H16]
Đoạn thẳng $AB$ dài $10\text{ cm}$. $M$ là trung điểm của đoạn thẳng $AB$. Độ dài đoạn thẳng $AM$ là
\choice
{$4\text{ cm}$}
{$6\text{ cm}$}
{$10\text{ cm}$}
{\True $5\text{ cm}$}
\loigiai{
$AM = 10 : 2 = 5\text{ cm}$.\\
\textbf{Đáp án của câu hỏi là D}
}
\end{ex}

\begin{ex} %[3H3H16]
Cho đoạn thẳng $CD$ có $I$ là trung điểm. Biết $CI = 6\text{ cm}$. Độ dài đoạn thẳng $CD$ là
\choice
{\True $12\text{ cm}$}
{$6\text{ cm}$}
{$18\text{ cm}$}
{$3\text{ cm}$}
\loigiai{
$CD = 6 \times 2 = 12\text{ cm}$.\\
\textbf{Đáp án của câu hỏi là A}
}
\end{ex}

\begin{ex} %[3H3H16]
Trên đoạn thẳng $MN = 16\text{ cm}$, điểm $K$ là trung điểm của $MN$. Độ dài đoạn thẳng $KN$ là
\choice
{$6\text{ cm}$}
{$10\text{ cm}$}
{$4\text{ cm}$}
{\True $8\text{ cm}$}
\loigiai{
$KN = 16 : 2 = 8\text{ cm}$.\\
\textbf{Đáp án của câu hỏi là D}
}
\end{ex}

\begin{ex} %[3H3H16]
Cho 3 điểm $P, Q, R$ thẳng hàng với $PQ = 4\text{ cm}, QR = 4\text{ cm}$. Khẳng định nào sau đây đúng?
\choice
{$P$ là trung điểm của $QR$}
{\True $Q$ là trung điểm của $PR$}
{$R$ là trung điểm của $PQ$}
{$Q$ không phải trung điểm của $PR$}
\loigiai{
$Q$ nằm giữa $P, R$ và $PQ = QR = 4\text{ cm}$ nên $Q$ là trung điểm của $PR$.\\
\textbf{Đáp án của câu hỏi là B}
}
\end{ex}

\begin{ex} %[3H3V16]
Đoạn thẳng $AB$ dài $24\text{ cm}$. $M$ là trung điểm của $AB$, $N$ là trung điểm của $AM$. Độ dài đoạn thẳng $AN$ là
\choice
{$4\text{ cm}$}
{$8\text{ cm}$}
{\True $6\text{ cm}$}
{$12\text{ cm}$}
\loigiai{
$AM = 24 : 2 = 12\text{ cm}$.\\
$AN = 12 : 2 = 6\text{ cm}$.\\
\textbf{Đáp án của câu hỏi là C}
}
\end{ex}

\begin{ex} %[3H3V16]
Một sợi dây dài $30\text{ cm}$. Bạn Mai gập đôi sợi dây lại để cắt lấy một nửa. Độ dài đoạn dây được cắt ra là
\choice
{$12\text{ cm}$}
{$20\text{ cm}$}
{$10\text{ cm}$}
{\True $15\text{ cm}$}
\loigiai{
$30 : 2 = 15\text{ cm}$.\\
\textbf{Đáp án của câu hỏi là D}
}
\end{ex}

\begin{ex} %[3H3V16]
Cho đoạn thẳng $AB$ dài $18\text{ cm}$. Điểm $M$ nằm giữa $A, B$ sao cho $AM = 9\text{ cm}$. Khẳng định nào sau đây đúng?
\choice
{$MB = 8\text{ cm}$}
{$M$ không phải trung điểm}
{\True $M$ là trung điểm của đoạn thẳng $AB$}
{$AB = 9\text{ cm}$}
\loigiai{
$MB = 18 - 9 = 9\text{ cm} = AM \Rightarrow M$ là trung điểm của $AB$.\\
\textbf{Đáp án của câu hỏi là C}
}
\end{ex}

\begin{ex} %[3H3V16]
Một khúc gỗ dài $40\text{ cm}$. Người ta cưa khúc gỗ đó tại chính giữa. Chiều dài mỗi đoạn gỗ sau khi cưa là
\choice
{$15\text{ cm}$}
{\True $20\text{ cm}$}
{$25\text{ cm}$}
{$30\text{ cm}$}
\loigiai{
$40 : 2 = 20\text{ cm}$.\\
\textbf{Đáp án của câu hỏi là B}
}
\end{ex}

\begin{ex} %[3H3V16]
Cho đoạn thẳng $MN$. Biết $O$ là trung điểm của $MN$ và $MO = 14\text{ cm}$. Độ dài đoạn thẳng $MN$ là
\choice
{$21\text{ cm}$}
{$32\text{ cm}$}
{\True $28\text{ cm}$}
{$24\text{ cm}$}
\loigiai{
$MN = 14 \times 2 = 28\text{ cm}$.\\
\textbf{Đáp án của câu hỏi là C}
}
\end{ex}

\begin{ex} %[3H3C16]
Cho đoạn thẳng $AB$ dài $32\text{ cm}$. Lấy điểm $M$ là trung điểm của $AB$, điểm $N$ là trung điểm của $MB$. Độ dài đoạn thẳng $AN$ là
\choice
{\True $24\text{ cm}$}
{$20\text{ cm}$}
{$26\text{ cm}$}
{$28\text{ cm}$}
\loigiai{
$AM = MB = 16\text{ cm}$.\\
$MN = NB = 16 : 2 = 8\text{ cm}$.\\
$AN = AM + MN = 16 + 8 = 24\text{ cm}$.\\
\textbf{Đáp án của câu hỏi là A}
}
\end{ex}

\textbf{Phần 2. Trắc nghiệm đúng sai}
\setcounterpageref{ex}{0}

\begin{ex} %[3H3H16]
Xét tính đúng sai về điểm ở giữa và trung điểm:
\choiceTF
{\True Điểm ở giữa hai điểm $A$ và $B$ phải thẳng hàng với $A$ và $B$.}
{\True Trung điểm của đoạn thẳng $AB$ chia đoạn thẳng $AB$ thành 2 đoạn bằng nhau.}
{Mọi điểm nằm ở giữa $A$ và $B$ đều là trung điểm của đoạn thẳng $AB$.}
{\True Chỉ có duy nhất một trung điểm của đoạn thẳng $AB$.}
\loigiai{
\begin{itemize}
    \item Điểm ở giữa phải thẳng hàng: Đúng.
    \item Trung điểm chia thành 2 phần bằng nhau: Đúng.
    \item Điểm ở giữa có thể cách không đều 2 đầu (chỉ trung điểm mới cách đều). Sai.
    \item Mỗi đoạn thẳng chỉ có 1 trung điểm duy nhất: Đúng.
\end{itemize}
\begin{itemize}
    \item a) \textbf{Đúng}.
    \item b) \textbf{Đúng}.
    \item c) \textbf{Sai}.
    \item d) \textbf{Đúng}.
\end{itemize}
}
\end{ex}

\begin{ex} %[3H3V16]
Cho đoạn thẳng $PQ = 20\text{ cm}$ có $M$ là trung điểm. Xét tính đúng sai:
\choiceTF
{\True Độ dài đoạn thẳng $PM$ là $10\text{ cm}$.}
{\True Độ dài đoạn thẳng $MQ$ là $10\text{ cm}$.}
{Nếu $K$ là trung điểm của $PM$ thì $PK = 6\text{ cm}$.}
{\True Nếu $K$ là trung điểm của $PM$ thì $PK = 5\text{ cm}$.}
\loigiai{
\begin{itemize}
    \item $PM = 20 : 2 = 10\text{ cm}$. Đúng.
    \item $MQ = 10\text{ cm}$. Đúng.
    \item $PK = 10 : 2 = 5\text{ cm} \neq 6\text{ cm}$. Ý 3 Sai, Ý 4 Đúng.
\end{itemize}
\begin{itemize}
    \item a) \textbf{Đúng}.
    \item b) \textbf{Đúng}.
    \item c) \textbf{Sai}.
    \item d) \textbf{Đúng}.
\end{itemize}
}
\end{ex}

\textbf{Phần 3. Trả lời ngắn}
\setcounterpageref{ex}{0}

\begin{ex} %[3H3H16]
Đoạn thẳng $AB$ dài $14\text{ cm}$. $M$ là trung điểm của $AB$. Tính độ dài đoạn thẳng $AM$ theo xăng-ti-mét.
\shortans{7}
\loigiai{
\textbf{Đáp số:} $7$\\
$14 : 2 = 7\text{ cm}$.
}
\end{ex}

\begin{ex} %[3H3H16]
Cho $I$ là trung điểm của đoạn thẳng $CD$. Biết $CI = 8\text{ cm}$. Tính độ dài đoạn thẳng $CD$ theo xăng-ti-mét.
\shortans{16}
\loigiai{
\textbf{Đáp số:} $16$\\
$8 \times 2 = 16\text{ cm}$.
}
\end{ex}

\begin{ex} %[3H3V16]
Đoạn thẳng $MN$ dài $40\text{ cm}$. $O$ là trung điểm của $MN$. Tính độ dài đoạn thẳng $ON$ theo xăng-ti-mét.
\shortans{20}
\loigiai{
\textbf{Đáp số:} $20$\\
$40 : 2 = 20\text{ cm}$.
}
\end{ex}

\begin{ex} %[3H3V16]
Cho đoạn thẳng $AB$ dài $36\text{ cm}$, $M$ là trung điểm của $AB$, $N$ là trung điểm của $AM$. Tính độ dài đoạn $AN$ theo xăng-ti-mét.
\shortans{9}
\loigiai{
\textbf{Đáp số:} $9$\\
$AM = 18\text{ cm} \Rightarrow AN = 18 : 2 = 9\text{ cm}$.
}
\end{ex}

\begin{ex} %[3H3V16]
Một thanh tre dài $50\text{ cm}$. Bác thợ cưa đúng ở trung điểm thanh tre thành 2 đoạn bằng nhau. Hỏi mỗi đoạn dài bao nhiêu xăng-ti-mét?
\shortans{25}
\loigiai{
\textbf{Đáp số:} $25$\\
$50 : 2 = 25\text{ cm}$.
}
\end{ex}

\begin{ex} %[3H3C16]
Cho đoạn thẳng $AB$ có trung điểm $M$. Biết $N$ là trung điểm của $AM$ và $AN = 5\text{ cm}$. Tính độ dài đoạn thẳng $AB$ theo xăng-ti-mét.
\shortans{20}
\loigiai{
\textbf{Đáp số:} $20$\\
$AM = 5 \times 2 = 10\text{ cm} \Rightarrow AB = 10 \times 2 = 20\text{ cm}$.
}
\end{ex}
